\begin{textsample}{POS Dim 1 – vem\_ed\_21 – Score 8.00 – vem\_ed\_21/t109.txt}  \label{ex:f1_pos_166}
Moderações Vozes emergentes, línguas e aprendizagem global Osvaldo Succi Junior fez a moderação de duas sessões de comunicação oral: “Leituras \textbf{colaborativas} e aprendizagem de línguas estrangeiras no projeto Literatandem”, por Vanessa Matiola (Unesp / Brasil) e “COIL en la enseñanza y el aprendizaje de \textbf{idiomas} extranjeros y segundas lenguas” (COIL no ensino e na aprendizagem de \textbf{idiomas} estrangeiros e segundas línguas), por Lilian Velasquez (DUOC UC / Colômbia) e Gabriel Cabezas (Universidad del Bío-Bío / Colômbia).
Ana Carolina Freschi fez a moderação do FlashTalk “Creative Global Learning at SHU: Focus on students’ engagement” (Aprendizagem Global Criativa na SHU: \textbf{Foco} no envolvimento dos alunos).
E apoiou Divinia Jithoo (Durban University of Technology / África do Sul) na moderação do painel “Emerging Voices” (Vozes Emergentes).
Estiveram presencialmente, no auditório principal do Centro de Convenções Rebouças, as vozes emergentes das painelistas Sandra Julieth Valencia Escobar (Universidad Católica De Manizales / Colômbia), Alia Gilbrecht (AN-Najah University / Palestina) e Bee Gan (Sheffield Hallam University / Reino Unido).
Daniel Otieno (Kenyatta University, Quênia) participou remotamente do painel.
Sandra Valencia \textbf{destacou} a \textbf{importância} de líderes de internacionalização nas instituições “sem os quais, o \textbf{que} ocorre \textbf{são} apenas iniciativas isoladas”.
E elencou os principais obstáculos à internacionalização sob a \textbf{perspectiva} latino-americana: barreira \textbf{linguística}, recursos econômicos, ferramentas tecnológicas e conectividade e \textbf{perspectiva} \textbf{política}.
Daniel Otieno ressaltou os \textbf{princípios} \textbf{que} devem guiar o design de projetos COIL: equidade, tomada de decisão democrática e translinguagem.
Alia Gilbrecht citou como principais desafios a captação de fundos, a conectividade e \textbf{questões} geopolíticas (como a guerra afeta a \textbf{todos}, física e emocionalmente).
Bee Gan sugere redesenhar os cursos para incluir a internacionalização no currículo.
Do \textbf{ponto} de vista estratégico, recomendou: “olhe para a missão da universidade, seus KPIs (Key Performance Indicators — indicadores-chave de desempenho) e mostre como os projetos COIL ajudam nisso.
Dessa \textbf{forma}, os Intercâmbios Virtuais ganham adesão \textbf{institucional}”.

% matched lemmas: colaborativo, destacar, foco, forma, idioma, importância, institucional, linguístico, perspectiva, político, ponto, princípio, que, questão, ser, todo
\end{textsample}
